% =====================================================================
%  統計基礎 4. 計算統計
%  構成は Docs/SectionPlan.md にしたがう。
% =====================================================================
\chapter{計算統計}
\label{chap:stat-computational}

\begin{chaptergoal}
  % TODO: この章の目標
\end{chaptergoal}

% ---------------------------------------------------------------------
\section{ブートストラップ}
\label{sec:stat-bootstrap}

\keywords{復元抽出、経験分布、リサンプリング}

% TODO: 動機：推定量の分布を式で求めるのが難しい場合がある

\subsection{経験分布}
% TODO:   各データに等しい確率 $1/n$ をおいた分布
% TODO:   経験分布関数と、真の分布関数との関係

\subsection{復元抽出}
% TODO:   取り出したものを戻してから次を取り出す抽出（非復元抽出との違い）

\subsection{リサンプリング}
% TODO:   手元のデータから復元抽出をくり返して「ブートストラップ標本」を作る
% TODO:   標準誤差の推定
% TODO:   パーセンタイル法による信頼区間
% TODO:   手順の整理と、Python による簡単な実装例
% TODO: 〔参考〕ジャックナイフ法

% ---------------------------------------------------------------------
\section{サンプリング}
\label{sec:stat-sampling}

\keywords{擬似乱数、逆変換法、棄却法、マルコフ連鎖モンテカルロ法}

\subsection{擬似乱数}
% TODO:   コンピュータで作る「乱数のように見える数」
% TODO:   線形合同法のしくみ、〔参考〕メルセンヌ・ツイスタ
% TODO:   シード（種）と再現性

\subsection{逆変換法}
% TODO:   一様乱数を累積分布関数の逆関数に入れると、目的の分布の乱数になる
% TODO:   例：指数分布の乱数
% TODO:   離散分布の場合

\subsection{棄却法}
% TODO:   作りやすい分布（提案分布）から乱数を作り、一定の確率で採用する
% TODO:   採用する確率の決め方と、なぜ正しい分布になるか
% TODO:   効率（採用される割合）

\subsection{マルコフ連鎖モンテカルロ法（MCMC）}
% TODO:   〔補足〕マルコフ連鎖：次の状態が今の状態だけで決まる確率的な動き、推移確率、定常分布
% TODO:   考え方：目的の分布を定常分布にもつマルコフ連鎖を作り、長く動かして乱数を得る
% TODO:   メトロポリス・ヘイスティングス法
% TODO:   ギブスサンプリング
% TODO:   バーンイン（初めの部分を捨てる）、収束の確かめ方
% TODO:   ベイズ統計の事後分布からのサンプリングへの応用

% ---------------------------------------------------------------------
\section{モンテカルロ積分}
\label{sec:stat-monte-carlo}

\keywords{モンテカルロ積分、期待値や確率密度の正規化定数、分散減少法}

\subsection{モンテカルロ積分}
% TODO:   積分を「期待値」として表し、乱数の平均で近似する
% TODO:   例：円周率 $\pi$ の近似
% TODO:   誤差は $1/\sqrt{n}$ の速さで小さくなる（中心極限定理による）
% TODO:   次元が高い積分でも使いやすいという利点

\subsection{期待値や確率密度の正規化定数}
% TODO:   期待値の近似計算
% TODO:   確率密度関数の正規化定数（全体を 1 にする定数）を乱数で近似する
% TODO:   ベイズ統計の周辺尤度の計算との関係

\subsection{分散減少法}
% TODO:   同じ乱数の数で、近似の誤差を小さくする工夫
% TODO:   重点サンプリング（重要な所から多く乱数をとる）
% TODO:   制御変量法
% TODO:   負の相関法（対称変量法）
% TODO:   層別サンプリング

% ---------------------------------------------------------------------
\section{欠測値の扱い}
\label{sec:stat-missing}

\keywords{EM 法}

% TODO: 〔補足〕欠測値の種類
% TODO:   完全にランダムな欠測（MCAR）、ランダムな欠測（MAR）、ランダムでない欠測（MNAR）
% TODO:   欠測のある行を捨てる方法（完全ケース分析）や平均で埋める方法の問題点

\subsection{EM 法}
% TODO:   考え方：欠けている部分を「期待値で補う（E ステップ）」と「パラメータを最尤推定する（M ステップ）」をくり返す
% TODO:   例1：一部が欠測した正規分布のデータ
% TODO:   例2：混合正規分布（どのグループから来たかが分からないデータ）
% TODO:   くり返すごとに尤度が減らないこと
% TODO:   注意：局所的な最大値に止まることがある、初期値の選び方
% TODO: 〔参考〕多重代入法

\begin{chaptersummary}
  % TODO: この章のまとめ
\end{chaptersummary}
